\documentclass[11pt]{ctexart}
\usepackage[a4paper,margin=1in]{geometry}
\usepackage{graphicx}
\usepackage{xcolor}
\usepackage{hyperref}
\definecolor{refgray}{gray}{0.45}
\title{\textbf{Market Views｜全球投行叙事汇编}\\[0.4em]{\large AI驱动存储与光互连结构性扩张，中国财政紧缩拖累GDP，人民币中间价进入关键观察期}}
\author{KC桌面 自动生成}
\date{260728}
\begin{document}
\maketitle
\tableofcontents
\newpage
\section{一页摘要}
\begin{itemize}
  \item AI基础设施从概念走向可量化订单，存储、光互连和内存接口芯片成为最清晰的结构性增长主线，订单出货比和积压订单覆盖倍数显示需求可见性显著改善。
  \item 全球存储市场供需缺口持续扩大，代理型AI驱动存储用量在2026-2030年增长超七倍，而供给端受洁净室、设备、材料等多重瓶颈制约，比特年复合增长率仅30-40\%。
  \item 内存接口芯片市场2030年规模预测存在分歧：Bernstein乐观情景约200亿美元（MRDIMM渗透率25-30\%），保守情景约80亿美元，渗透率节奏是核心变量。
  \item 中国二季度GDP增速放缓至4.3\%，低于全年目标下沿，主要受广义财政赤字率大幅收窄影响，而非需求不足。
  \item 中国互联网平台反垄断处罚落地消除不确定性，但消费品牌面临传统渠道调整压力，旺旺传统批发渠道收入双位数下滑。
  \item Oticon Zeal在VA渠道的份额爬坡速度（两个月+420bp）远超行业历史新品，标志助听器市场进入技术驱动的新竞争阶段。
  \item NOM模型显示人民币中间价预测值下移254个基点至6.7685，主要受美元走弱驱动，而非中国央行主动引导；逆周期因子计入后仍偏强，后续政策会议密集。
  \item 中国在全球资本品出口份额持续上升（2005年7\%→2025年24\%），欧洲份额相应下降（54\%→43\%），对非美发达市场商品价格形成持续下行压力。
\end{itemize}
\section{全球投行叙事汇编}
今日纳入 32 篇投行/券商研究，按 5 个市场、资产与产业主题系统整理。
\newpage
\subsection{AI驱动存储与光互连结构性扩张，半导体周期进入需求验证阶段}
\textbf{AI基础设施正从概念走向可量化的订单与收入，存储、光互连和内存接口芯片成为最清晰的结构性增长主线；同时，汽车半导体在库存过低与中国电动车需求推动下启动复苏，但大模型参数竞赛的边际回报递减，效率与系统集成成为新竞争维度。}
\subsubsection*{主流共识}
\begin{itemize}
  \item AI数据中心对存储、光互连和功率半导体的需求正在转化为可量化的订单和收入，订单出货比和积压订单覆盖倍数均显示需求可见性显著改善。
  \item 全球存储市场供需缺口持续扩大，代理型AI驱动存储用量在2026-2030年增长超七倍，而供给端受洁净室、设备、材料等多重瓶颈制约，比特年复合增长率仅30-40\%。
  \item 汽车半导体库存已降至不可持续的低水平，叠加中国电动车需求，复苏具有持续性而非一次性补库。
\end{itemize}
\subsubsection*{分歧与条件差异}
\begin{itemize}
  \item 内存接口芯片市场2030年规模预测存在分歧：Bernstein乐观情景下约200亿美元（MRDIMM渗透率25-30\%），保守情景下约80亿美元，渗透率节奏是核心变量。
  \item 中国大模型竞争焦点从参数规模转向工程效率与系统集成，部分机构认为3万亿参数是下一个里程碑，但专家指出纯参数扩展边际回报递减，模型-智能体集成系统成为新壁垒。
  \item 英伟达入股Naver的模式是否可复制：Bernstein认为这标志着AI基础设施融资从单一资产负债表转向伙伴驱动的“新云”结构，但其他投行尚未明确表态。
\end{itemize}
\subsubsection*{机构观点}
\paragraph{Bernstein} 全球内存接口芯片市场2030年规模上调至约200亿美元（65\% CAGR），MRDIMM升级使单模组接口硅含量达传统RDIMM的10倍，MRCD和MDB芯片将贡献73\%市场。三家头部供应商（澜起、瑞萨、Rambus）合计份额超90\%，Rambus受益于代理型AI驱动的服务器CPU复兴、DRAM专利许可和AI ASIC IP授权三大引擎，但历史诉讼限制其份额增长，研发效率低于澜起。
{\small\color{refgray}数据：全球内存接口芯片市场2030年规模约200亿美元，对应65\% CAGR\par}
{\small\color{refgray}数据：MRDIMM模组接口硅含量为传统RDIMM的10倍\par}
{\small\color{refgray}数据：MRCD和MDB芯片贡献2030年市场73\%规模\par}
{\small\color{refgray}数据：三家头部供应商2024年合计份额超90\%\par}
{\small\color{refgray}边际变化：Bernstein上调TAM预测，核心假设MRDIMM渗透率2030年达25-30\%，此前市场预期更为保守。\par}
\paragraph{Bernstein} 韩国AI峰会上，SK海力士与英伟达签署超5000亿美元存储合作意向书，三星与博通签署2000亿美元MOU（含HBM与2nm代工），SK集团另有2500亿美元额外合作。这些框架性协议显示英伟达和博通正用长期协议锁定未来数年存储供应，尤其是HBM4及后续世代。全球存储行业年收入2027-2028年预计约1.3万亿美元，上述合作规模占相当份额，有助于缓解供应能力担忧。
{\small\color{refgray}数据：SK海力士与英伟达合作规模超5000亿美元\par}
{\small\color{refgray}数据：三星与博通MOU价值2000亿美元（至2030年）\par}
{\small\color{refgray}数据：SK集团另有2500亿美元额外合作\par}
{\small\color{refgray}数据：全球存储年收入2027-2028年约1.3万亿美元\par}
{\small\color{refgray}边际变化：存储合作从概念性框架转向具体金额和期限，英伟达和博通首次以如此大规模长期协议锁定存储供应。\par}
\paragraph{Bernstein} 英伟达以10亿美元认购Naver约5\%股份，成为战略股东。这一动作标志着AI基础设施融资模式从单一资产负债表转向伙伴驱动的“新云”结构：基础设施投资方出资，芯片供应商支持算力部署，模型方提供智能。英伟达入股直接降低了市场对Naver AI工厂执行和融资能力的疑虑，后续SPV融资和外部资本引入预计更为顺畅。
{\small\color{refgray}数据：英伟达以10亿美元认购Naver约720万股，持股约5\%\par}
{\small\color{refgray}数据：发行价每股204,500韩元\par}
{\small\color{refgray}数据：Naver、英伟达和Brookfield计划共同研究千兆瓦级AI云基础设施\par}
{\small\color{refgray}边际变化：英伟达从芯片供应商转变为平台股东，利益绑定从客户关系升级为股权关系，为AI基础设施融资提供新模式。\par}
\paragraph{GS} 意法半导体数据中心收入指引大幅上调：2026年超10亿美元，2027年远高于20亿美元，主要驱动力来自800G和1.6T光模块。订单出货比接近2倍，光互连领域超2倍，积压订单覆盖4.5-5个季度收入，2026年数据中心收入已被订单完全覆盖。毛利率改善路径清晰，但200mm→300mm硅生产迁移带来约60bp影响，预计持续至2027年底。
{\small\color{refgray}数据：2026年数据中心收入超10亿美元，2027年远高于20亿美元\par}
{\small\color{refgray}数据：整体订单出货比接近2倍，光互连超2倍\par}
{\small\color{refgray}数据：积压订单相当于4.5-5个季度收入\par}
{\small\color{refgray}数据：制程迁移带来约60bp毛利率影响至2027年底\par}
{\small\color{refgray}边际变化：GS大幅上调数据中心收入预期，且订单覆盖倍数显示需求从改善走向可量化，此前市场仅预期温和增长。\par}
\paragraph{GS} 诺基亚Q2光学订单从Q1的10亿欧元跃升至28亿欧元，AI和云客户成为主要增长引擎，管理层预计约半数订单将在12个月内确认为收入。磷化铟晶圆和存储器等前沿组件供应持续紧张，客户通过更长期订单锁定产能。诺基亚正通过收购NXP亚利桑那工厂、圣何塞工厂量产、宾夕法尼亚工厂扩产（测试封装能力提升10倍）来应对瓶颈。
{\small\color{refgray}数据：Q2光学订单28亿欧元，Q1为10亿欧元\par}
{\small\color{refgray}数据：约半数订单将在12个月内确认为收入\par}
{\small\color{refgray}数据：宾夕法尼亚工厂光学测试封装能力提升约10倍\par}
{\small\color{refgray}数据：2026年预计产生约8亿欧元重组费用\par}
{\small\color{refgray}边际变化：光学订单季度环比增长近3倍，且收入确认节奏快于市场预期，此前市场认为订单转化需更长时间。\par}
\paragraph{MS} 汽车半导体复苏由库存过低和中国电动车需求双驱动：TI确认汽车客户库存降至不可持续低水平，STMicro汽车业务Q2同比增长13-14\%，SiC收入环比增长超30\%且book-to-bill远高于1。NXP的结构性优势来自软件定义汽车架构转型（S32N、成像雷达等），加速增长驱动已占汽车收入45\%以上。数据中心对功率半导体的拉动正从概念走向可量化收入。
{\small\color{refgray}数据：STMicro汽车业务Q2同比增长约13-14\%\par}
{\small\color{refgray}数据：STMicro SiC收入环比增长超30\%\par}
{\small\color{refgray}数据：NXP加速增长驱动占汽车收入45\%以上\par}
{\small\color{refgray}数据：TI汽车客户库存降至不可持续低水平\par}
{\small\color{refgray}边际变化：汽车半导体复苏被确认为有持续性的周期启动，而非一次性补库，此前市场担忧需求可持续性。\par}
\paragraph{NOM} 全球存储供需缺口短期难解：代理型AI驱动存储用量2026-2030年增长超七倍（即使考虑四倍内存效率折扣），而供给端比特年复合增长率仅30-40\%。长鑫存储产能扩张节奏高于行业均值（40-45\% CAGR），全球DRAM市占率有望从约10\%提升至2028年底约18\%。财务预测2026-2028年营收CAGR 63\%，归母净利润CAGR 74\%，毛利率从41\%提升至90\%以上。
{\small\color{refgray}数据：全球存储用量2026-2030年增长超七倍\par}
{\small\color{refgray}数据：长鑫存储比特扩张CAGR 40-45\%\par}
{\small\color{refgray}数据：全球DRAM市占率从约10\%提升至2028年底约18\%\par}
{\small\color{refgray}数据：2026-2028年营收CAGR 63\%，归母净利润CAGR 74\%\par}
{\small\color{refgray}边际变化：NOM首次覆盖长鑫存储，给出积极展望，此前市场对中国DRAM厂商扩张能力存在疑虑。\par}
\paragraph{NOM} WinSemi光学代工业务从样品阶段转向量产爬坡：光电探测器（PD）已为单一战略客户量产，产能接近满载；VCSEL和长距离EML已量产，中距离EML预计2H26加入，CW激光器预计2027年下半年贡献收入。2026年光学通信收入预计占高个位数百分比，2027年有望迈向双位数。毛利率改善来自产品组合优化（基础设施收入占比已与手机相当）和折旧接近底部。
{\small\color{refgray}数据：Q2毛利率28.2\%，营业利润率14.1\%\par}
{\small\color{refgray}数据：产能利用率从60\%提升至65\%\par}
{\small\color{refgray}数据：折旧费用每季约7.6亿新台币，接近底部\par}
{\small\color{refgray}数据：Q3毛利率指引升至低30\%区间\par}
{\small\color{refgray}边际变化：光学代工从样品阶段进入量产爬坡，此前市场仅预期样品阶段贡献有限。\par}
\paragraph{NOM} 中国大模型下一阶段竞争焦点从参数规模转向工程效率与系统集成。约3万亿参数是下一个里程碑，但纯参数扩展边际回报递减，模型竞争力取决于架构效率、训练后优化、内存管理和长周期任务表现。DeepSeek推迟V4发布部分原因是在智能体框架层面调试，目标构建类似Claude Code的一体化产品，创造专有数据飞轮。模型代际领先时间已缩短至3-4个月。
{\small\color{refgray}数据：约3万亿参数是中国大模型下一个里程碑\par}
{\small\color{refgray}数据：DeepSeek推理成本约0.8元/百万token（60\%缓存命中率）\par}
{\small\color{refgray}数据：推理毛利率约40\%\par}
{\small\color{refgray}数据：模型代际领先时间缩短至3-4个月\par}
{\small\color{refgray}边际变化：竞争焦点从参数规模转向效率与系统集成，此前市场主要关注参数规模排名。\par}
\subsubsection*{关键数据}
\begin{itemize}
  \item 全球内存接口芯片市场2030年规模约200亿美元（Bernstein）
  \item SK海力士与英伟达合作规模超5000亿美元（Bernstein）
  \item 全球存储年收入2027-2028年约1.3万亿美元（Bernstein）
  \item 意法半导体2026年数据中心收入超10亿美元，2027年远高于20亿美元（GS）
  \item 诺基亚Q2光学订单28亿欧元，环比增长近3倍（GS）
  \item 长鑫存储全球DRAM市占率有望从10\%提升至2028年底约18\%（NOM）
  \item STMicro汽车业务Q2同比增长13-14\%，SiC收入环比增长超30\%（MS）
  \item 中国大模型约3万亿参数是下一个里程碑，但纯参数扩展边际回报递减（NOM）
\end{itemize}
\begin{figure}[htbp]
\centering
\includegraphics[width=0.88\linewidth]{figures/fig_007.jpg}
\caption{EXHIBIT 2 - EXHIBIT 3: We raise our global memory interface TAM projection to USD 20Bn in 2030, with MRDIMM's MRCD and MDB contributing 73\%, driven by significantly higher chipset value per module Global memory interface chips TAM (Bn USD)}
\end{figure}
\begin{figure}[htbp]
\centering
\includegraphics[width=0.88\linewidth]{figures/fig_011.jpg}
\caption{EXHIBIT 1 - EXHIBIT 1: Samsung 2.5D Cube-S is similar to TSMC CoWoS-S, and 2.3D Cube-E \& -R are similar to TSMC CoWoS-L \& -R, respectively. EXHIBIT 2: TSMC offers three variants of CoWoS. EXHIBIT 3: Consensus memory industry to generate c.}
\end{figure}
\begin{figure}[htbp]
\centering
\includegraphics[width=0.88\linewidth]{figures/fig_006.jpg}
\caption{EXHIBIT 1 - EXHIBIT 1: Rambus has two primary business segments but reports financial statements in three lines. Two business segments ride on the wave of CPU Renaissance and AI ASIC respectively The revenue mix and growth rate numbers are o}
\end{figure}
\begin{figure}[htbp]
\centering
\includegraphics[width=0.88\linewidth]{figures/fig_057.jpg}
\caption{Exhibit 8 - Exhibit 8: Weekly stock performance Exhibit 9: Monthly stock performance Exhibit 10: YTD stock performance}
\end{figure}
{\small\color{refgray}本节综合 9 篇研究；涉及 Bernstein、GS、MS、NOM。}
\newpage
\subsection{消费与互联网：监管落地与结构性分化}
\textbf{中国互联网平台反垄断处罚落地消除不确定性，但消费品牌面临渠道与需求的结构性分化，奢侈品旅游消费亦因客群依赖度差异而表现不一。}
\subsubsection*{主流共识}
\begin{itemize}
  \item 携程反垄断罚款金额高于阿里和美团，但市场更关注业务调整的长期影响而非一次性财务支出。
  \item 奢侈品牌欧洲市场表现分化，对中国游客依赖度低、价差小的品牌表现更稳健。
  \item 中国消费品牌面临传统渠道调整压力，旺旺传统批发渠道收入双位数下滑。
\end{itemize}
\subsubsection*{分歧与条件差异}
\begin{itemize}
  \item 携程排他协议取消后，美团和抖音等竞争对手能否缩小酒店供给差距存在分歧：GS认为监管未额外限制市场份额，NOM则预期竞争加剧。
  \item 奢侈品牌中，美国游客消费能否有效缓冲中国游客减少的影响，因品牌而异：历峰受益于美国需求，蒙克莱则缺乏此缓冲。
  \item 印度消费公司Tata Consumer的增长业务（47\%增速）能否持续高增长，取决于新品推出和品类拓展，但基数效应可能放缓。
\end{itemize}
\subsubsection*{机构观点}
\paragraph{GS} 携程反垄断罚款51.79亿元（基于2025年国内收入7.5\%），高于阿里（4\%）和美团（3\%），但公司持有1040亿元流动性，足以覆盖。监管未对抽佣率、市场份额或与同程合作施加额外限制，市场可能正面解读。
{\small\color{refgray}数据：罚款51.79亿元，占2025年国内收入7.5\%\par}
{\small\color{refgray}数据：携程持有1040亿元流动性\par}
{\small\color{refgray}数据：阿里和美团处罚后股价分别涨8\%和9\%\par}
{\small\color{refgray}边际变化：监管终局信号比罚款金额更关键，此前市场担忧的不确定性消除。\par}
\paragraph{GS} 旺旺中国FY1Q26收入同比-6\%，净利润同比-38\%，低于市场预期。传统批发渠道收入双位数下滑，运营费用高个位数增长。当前PE约9倍，低于可比公司卫龙（11倍）、洽洽（13倍）、盐津铺子（14倍），反映市场对其增长路径的怀疑。
{\small\color{refgray}数据：FY1Q26收入同比-6\%，净利润同比-38\%\par}
{\small\color{refgray}数据：传统批发渠道收入双位数下滑\par}
{\small\color{refgray}数据：运营费用同比高个位数增长\par}
{\small\color{refgray}数据：PE约9倍，可比公司11-14倍\par}
{\small\color{refgray}边际变化：公司处于旧渠道失效、新投入尚未见效的过渡期，估值折价隐含市场对过渡期延长的预期。\par}
\paragraph{NOM} 携程被要求停止排他协议和最低价要求，罚款约53亿元占净现金11\%。取消排他后，美团和抖音可能缩小酒店供给差距，参考2021年案例，竞争对手份额从17\%涨至40\%。预计FY26/FY27非GAAP经营利润率分别下降3pp/1pp。
{\small\color{refgray}数据：罚款约53亿元，占净现金11\%\par}
{\small\color{refgray}数据：2025年在线酒店预订市场份额近60\%\par}
{\small\color{refgray}数据：参考案例：竞争对手份额从17\%涨至40\%\par}
{\small\color{refgray}数据：预计FY26/FY27经营利润率下降3pp/1pp\par}
{\small\color{refgray}边际变化：排他协议取消后，酒店产品在各平台间技术性同质化，竞争加剧，但携程一站式模式仍维持最大份额。\par}
\paragraph{NOM} Tata Consumer印度业务收入同比+12\%，营业利润+19.3\%。增长业务（即饮、Sampann等）占比升至36\%，同比增速47\%，成为主要驱动力。核心茶盐业务稳健，茶销量+2\%，盐销量+7\%。管理层预期全年OPM扩张50-75bp。
{\small\color{refgray}数据：印度业务收入同比+12\%，营业利润+19.3\%\par}
{\small\color{refgray}数据：增长业务占比36\%，同比增速47\%\par}
{\small\color{refgray}数据：茶销量+2\%，盐销量+7\%\par}
{\small\color{refgray}数据：毛利率同比+255bp至42.7\%\par}
{\small\color{refgray}边际变化：增长业务以47\%增速重塑公司DNA，核心业务通过提价（茶+1-2\%，盐+6\%）支撑利润率，结构优化加速。\par}
\paragraph{Bernstein} 奢侈品牌欧洲市场表现分化由三个因素决定：对中国游客依赖度、区域价差、美国游客替代能力。硬奢（历峰）中国客群占比20-25\%，价差小，欧洲区域有机增长+7\%；蒙克莱中国客群占比35-40\%，价差30\%，欧洲区域有机收入-8\%。
{\small\color{refgray}数据：历峰中国客群占比20-25\%，欧洲区域有机增长+7\%\par}
{\small\color{refgray}数据：蒙克莱中国客群占比35-40\%，中法价差30\%\par}
{\small\color{refgray}数据：蒙克莱欧洲区域有机收入-8\%\par}
{\small\color{refgray}数据：软奢品牌平均价差约20\%\par}
{\small\color{refgray}边际变化：美国游客消费成为部分品牌缓冲变量，但蒙克莱缺乏此调节空间，季节性特征放大。\par}
\subsubsection*{关键数据}
\begin{itemize}
  \item 携程罚款51.79亿元，占2025年国内收入7.5\%
  \item 旺旺FY1Q26净利润同比-38\%，PE约9倍
  \item 历峰欧洲区域有机增长+7\%，蒙克莱-8\%
  \item Tata Consumer增长业务占比36\%，同比增速47\%
  \item 携程持有1040亿元流动性
  \item 蒙克莱中法价差30\%，软奢平均20\%
\end{itemize}
\begin{figure}[htbp]
\centering
\includegraphics[width=0.88\linewidth]{figures/fig_014.jpg}
\caption{EXHIBIT 1 - FIGURE 1: Shopping in Europe has remained relatively attractive for Chinese consumers looking to purchase Hermès and Moncler products - even after accounting for near-shore alternatives such as Japan. More recent cross-continental}
\end{figure}
\begin{figure}[htbp]
\centering
\includegraphics[width=0.88\linewidth]{figures/fig_034.jpg}
\caption{Exhibit 1 - Exhibit 1: TCOM earnings breakdown Exhibit 2: Alibaba and Meituan share price during and after SAMR investigations Exhibit 4: TCOM historical forward P/E ex. exceptional}
\end{figure}
\begin{figure}[htbp]
\centering
\includegraphics[width=0.88\linewidth]{figures/fig_015.jpg}
\caption{FIGURE 1 - EXHIBIT 3: Strong demand in the USA, on the other hand, presents a potential offset for luxury brands come 2Q26E, if some of this overflows into Europe as a sequential improvement in American tourist spending. Moncler, faced with r}
\end{figure}
\begin{figure}[htbp]
\centering
\includegraphics[width=0.88\linewidth]{figures/fig_036.jpg}
\caption{Exhibit 3 - Exhibit 4: TCOM historical forward P/E ex. exceptional \#\# Price Target Risks and Methodology – Trip.com Group Valuation methodology: We have a Buy rating on Trip.com (TCOM/9961.HK) with 12-month target prices of US\textbackslash{}\$71/HK\textbackslash{}\$560 ba}
\end{figure}
{\small\color{refgray}本节综合 5 篇研究；涉及 Bernstein、GS、NOM。}
\newpage
\subsection{医疗与工业：结构性分化与数据驱动拐点}
\textbf{医疗板块聚焦助听器与肿瘤疫苗的竞争格局重塑及关键数据催化，工业板块关注中国资本品出口对欧洲的份额侵蚀与造船业估值修复。}
\subsubsection*{主流共识}
\begin{itemize}
  \item Oticon Zeal在VA渠道的份额爬坡速度（两个月+420bp）远超行业历史新品，标志助听器市场进入技术驱动的新竞争阶段。
  \item Moderna的黑色素瘤三期中期数据是肿瘤疫苗平台价值的关键验证节点，市场已部分定价成功预期。
  \item 中国在全球资本品出口份额持续上升（2005年7\%→2025年24\%），欧洲份额相应下降（54\%→43\%），且追赶在优势领域加速。
  \item 三星重工估值回调后进入合理区间，但中国船厂产能扩张对新造船价格形成持续压力。
\end{itemize}
\subsubsection*{分歧与条件差异}
\begin{itemize}
  \item Oticon Zeal的份额增长是否可持续：GS认为趋势刚启动，但需关注定制式助听器占比提升对品牌格局的长期影响。
  \item Moderna肿瘤疫苗成功概率：期权市场隐含成功时上涨49\% vs 失败时回归基础估值约\$25，分歧在于市场是否过度乐观。
  \item Vertex的SION CF数据影响：市场以汗液氯离子降低10mmol/L为门槛，但GS认为Vertex内部标准更高，需结合产品组合判断。
  \item 中国对欧出口加速是结构性还是周期性：GS指出双边贸易总量扩大，但“出口增长+进口停滞”品类是竞争力转移的早期信号。
\end{itemize}
\subsubsection*{机构观点}
\paragraph{GS} Oticon Zeal在VA渠道上市两个月内价值份额从18.9\%升至23.1\%（+420bp），爬坡速度是前代Intent（+210bp）的两倍，远超Sonova和GN同期新品（+70bp和+60bp）。6月VA渠道整体销量同比+5\%，销售额+13.8\%，定制式助听器价值占比从15\%升至22\%，反映行业从“量增”转向“价增”。Oticon销量同比+27\%，销售额+55\%，预计对Demant Q2有机增长贡献240bp。
{\small\color{refgray}数据：Oticon VA份额从18.9\%升至23.1\%（两个月+420bp）\par}
{\small\color{refgray}数据：前代Intent同期仅+210bp\par}
{\small\color{refgray}数据：VA渠道6月销量+5\%，销售额+13.8\%\par}
{\small\color{refgray}数据：定制式助听器价值占比从15\%升至22\%\par}
{\small\color{refgray}边际变化：Oticon Zeal的份额爬坡速度显著快于历史新品，定制式助听器占比快速提升，行业增长逻辑从“卖更多”转向“卖更贵”。\par}
\paragraph{GS} Moderna黑色素瘤辅助治疗三期中期数据是肿瘤板块价值的关键节点，当前股价已部分反映成功预期。期权市场显示成功时到明年1月上涨49\%，失败则回归基础业务估值约\$25。公司Q2收入指引5000万-1亿美元（GS预测9000万，共识1.03亿），FY26指引同比增速≤10\%（GS/共识7.2\%/7.8\%），已考虑美国新冠疫苗接种量下降。
{\small\color{refgray}数据：Moderna年初至今跑赢生物科技板块61\%\par}
{\small\color{refgray}数据：期权市场隐含成功时上涨49\%，失败估值约\$25\par}
{\small\color{refgray}数据：Q2收入指引5000万-1亿美元\par}
{\small\color{refgray}数据：FY26指引同比增速≤10\%\par}
{\small\color{refgray}边际变化：市场对黑色素瘤数据已部分定价，但期权隐含波动揭示事件二元性；收入指引稳定但增长动力尚未恢复。\par}
\paragraph{GS} Vertex Q2收入预计31.55亿美元（共识32.19亿），非GAAP EPS 4.57美元（共识4.70）。下半年核心催化剂：povetacicept（IgA肾病）11月30日PDUFA，现有血尿改善和每月给药方案支持同类最佳潜力；SION CF数据（8月）是短期压制，但Vertex内部标准更高；收购Crinetics预计Q3完成，atumelnant（CAH/Cushing）二期数据待读出。
{\small\color{refgray}数据：Q2收入预计31.55亿美元\par}
{\small\color{refgray}数据：povetacicept PDUFA日期11月30日\par}
{\small\color{refgray}数据：SION CF二期数据8月公布\par}
{\small\color{refgray}数据：收购Crinetics预计Q3完成\par}
{\small\color{refgray}边际变化：市场焦点从Q2业绩转向下半年密集数据催化，povetacicept上市潜力是中期价值锚点，SION数据是短期压制。\par}
\paragraph{GS} 中国在全球资本品出口份额从2005年7\%升至2025年24\%，欧洲从54\%降至43\%。截至2026年6月，中国对欧盟出口增速从7.6\%跃升至18.5\%，43个品类三个月滚动同比增速从9.6\%加速至15.4\%。增速最快品类：钨制品（+233\%）、光纤电缆（+111\%）、船用发动机（+75\%）、热泵（+60\%）、轻型商用车（+60\%）。欧洲在商用车发动机、乳品机械、铁路制动器仍具相对优势，但过去五年中国在这些领域份额增长最快。
{\small\color{refgray}数据：中国资本品出口份额从7\%升至24\%\par}
{\small\color{refgray}数据：欧洲份额从54\%降至43\%\par}
{\small\color{refgray}数据：中国对欧盟出口增速从7.6\%升至18.5\%\par}
{\small\color{refgray}数据：钨制品出口增速+233\%\par}
{\small\color{refgray}边际变化：中国对欧出口增速显著加快，且在欧洲优势领域追赶速度最快，贸易防御调查数量激增至10年均值2倍以上。\par}
\paragraph{NOM} 三星重工Q2营业利润3250亿韩元（同比+58.7\%），低于共识13.2\%，受250亿韩元退休金一次性成本影响。新船坞贡献1500亿韩元额外收入，但运营利润率低于造船业务均值。2026年新订单预计170亿美元（同比+115.7\%），F-LNG是主要驱动力（预计80亿美元，同比+900\%）。当前P/B 3.7倍，目标价25000韩元（隐含+9.2\%），估值已进入合理区间。
{\small\color{refgray}数据：Q2营业利润3250亿韩元，低于共识13.2\%\par}
{\small\color{refgray}数据：一次性成本250亿韩元\par}
{\small\color{refgray}数据：新船坞收入1500亿韩元，利润率低于均值\par}
{\small\color{refgray}数据：2026年新订单预计170亿美元（+115.7\%）\par}
{\small\color{refgray}边际变化：报价自4月高点回调33.4\%后估值进入合理区间，但新产能利润率偏低限制评级上调空间，F-LNG订单是主要增量。\par}
\subsubsection*{关键数据}
\begin{itemize}
  \item Oticon VA份额两个月+420bp，前代仅+210bp
  \item VA渠道定制式助听器价值占比从15\%升至22\%
  \item Moderna期权隐含成功时上涨49\%，失败估值约\$25
  \item Vertex povetacicept PDUFA 11月30日
  \item 中国资本品出口份额从7\%升至24\%，欧洲从54\%降至43\%
  \item 中国对欧盟出口增速从7.6\%升至18.5\%
  \item 三星重工Q2营业利润低于共识13.2\%
  \item 三星重工2026年F-LNG新订单预计80亿美元（+900\%）
\end{itemize}
\begin{figure}[htbp]
\centering
\includegraphics[width=0.88\linewidth]{figures/fig_017.jpg}
\caption{Exhibit 1 - Exhibit 2: Oticon saw strong y/y value growth Oticon y/y value growth in the VA channel Exhibit 3: ...and reached all-time highs in VA value share in June 2026 (2nd month since its the launch in the VA channel in May) Oticon - Va}
\end{figure}
\begin{figure}[htbp]
\centering
\includegraphics[width=0.88\linewidth]{figures/fig_024.jpg}
\caption{Exhibit 1 - Exhibit 1: LED lighting, excavators and heat pumps have shown the most growth in Chinese exports globally (ex Russia) and greatest slowdown in imports over the last decade Snapshot of competitive intensity - China Size of bubbl}
\end{figure}
\begin{figure}[htbp]
\centering
\includegraphics[width=0.88\linewidth]{figures/fig_025.jpg}
\caption{Exhibit 1 - Exhibit 3: Engines, dairy machinery and rail brakes are the areas where Europe is most dominant vs. China in global exports markets, but the gap has been decreasing Snapshot of competitive intensity - Global, Europe/China market sh}
\end{figure}
\begin{figure}[htbp]
\centering
\includegraphics[width=0.88\linewidth]{figures/fig_027.jpg}
\caption{Exhibit 4 - Exhibit 4: China net exports to Europe have expanded significantly and regained the rising trend post covid highs China's trade balance (net export) with Europe for 43 cap goods categories, million \textbackslash{}\$ - 12m rolling ■ China Trade}
\end{figure}
{\small\color{refgray}本节综合 5 篇研究；涉及 GS、NOM。}
\newpage
\subsection{宏观与市场主题：中国财政紧缩、AI算力需求与全球通胀新变量}
\textbf{中国广义财政赤字率大幅收窄是二季度GDP放缓的主因，但AI算力需求的结构性增长和全球贸易格局变化正在重塑通胀与产业投资逻辑。}
\subsubsection*{主流共识}
\begin{itemize}
  \item 中国二季度GDP增速放缓至4.3\%，低于全年目标下沿，主要受财政紧缩影响，而非需求不足。
  \item AI算力需求长期供不应求，近期板块回调属技术性调整，基本面未变。
  \item 中国贸易模式转变（多出少进）对非美发达市场商品价格形成持续下行压力。
\end{itemize}
\subsubsection*{分歧与条件差异}
\begin{itemize}
  \item 美联储7月利率路径：通胀改善支持按兵不动，但地缘政治风险推升加息概率至35\%。
  \item 中国数据中心需求：上半年芯片供应导致节奏放缓，但全年需求预测仍维持20\%复合增长率。
  \item HDD行业定价：市场担忧产能过剩，但机构认为结构性供应缺口将持续至2028年，合约价已上移至25-30美元/TB。
\end{itemize}
\subsubsection*{机构观点}
\paragraph{GS} 中国二季度GDP放缓至4.3\%的主因是广义财政赤字率从2025年的11.0\%收窄至7.0\%，土地出让收入6月同比骤降42\%是核心拖累。7月政治局会议预计释放更强宽松信号，但地方财政约束下执行能力仍是关键。
{\small\color{refgray}数据：6月预算内财政支出增速4.0\%，远低于收入增速8.7\%\par}
{\small\color{refgray}数据：广义财政赤字率从11.0\%收窄至7.0\%\par}
{\small\color{refgray}数据：土地出让收入6月同比-42\%\par}
{\small\color{refgray}边际变化：此前市场更多关注需求端疲弱，GS明确将财政紧缩而非需求不足作为二季度放缓的主因，并强调土地出让收入持续下滑对地方财力的系统性削弱。\par}
\paragraph{GS} 中国数据中心上半年智能算力新增约595 EFLOPS，低于2H25的802 EFLOPS，主因芯片供应瓶颈。但需求逻辑未变：AI开源模型推动token扩张，智谱已启动1GW数据中心建设。维持2025-2028年20\%复合增长率，总容量从16GW增至27GW。
{\small\color{refgray}数据：1H26新增智能算力595 EFLOPS，2H25为802 EFLOPS\par}
{\small\color{refgray}数据：2026年并网需求预测下调至约3GW\par}
{\small\color{refgray}数据：2025-2028年CAGR 20\%，总容量从16GW增至27GW\par}
{\small\color{refgray}边际变化：此前市场担忧需求走弱，GS指出节奏放缓是供给侧阶段性约束，而非需求转向，并强调国产硬件（华为Atlas 950）性能差距从1/8缩小至1/4。\par}
\paragraph{GS} 中国贸易模式转变（出口增长约20\%，进口低于疫情前趋势）已使非美发达市场商品价格平均下降0.7\%，日本和欧元区分别达1.1\%和1.0\%。出口渗透率每提升1个百分点，商品价格下降0.5\%。
{\small\color{refgray}数据：中国对非美发达市场出口增长约20\%\par}
{\small\color{refgray}数据：出口渗透率每升1pp，商品价格降0.5\%\par}
{\small\color{refgray}数据：非美发达市场商品价格平均降0.7\%\par}
{\small\color{refgray}数据：日本降1.1\%，欧元区降1.0\%\par}
{\small\color{refgray}边际变化：此前市场关注中国出口对全球通胀的推升作用，GS量化了出口增长与进口减弱共同带来的通缩效应，并指出这一影响将持续。\par}
\paragraph{GS} 6月核心CPI环比下降2个基点，通胀改善支持美联储按兵不动，但地缘政治风险（伊朗局势、俄罗斯炼油厂袭击）推高能源价格，加息概率从25\%上调至35\%。市场定价隐含约40\%加息概率，若维持将成近几十年最大意外之一。
{\small\color{refgray}数据：6月核心CPI环比-0.02\%\par}
{\small\color{refgray}数据：加息概率从25\%上调至35\%\par}
{\small\color{refgray}数据：市场定价隐含约40\%加息概率\par}
{\small\color{refgray}边际变化：此前GS基准预测为年底前按兵不动，但地缘政治风险上升使加息辩论重新活跃，且冲突对加息的影响可能超过油价传导的数学计算。\par}
\paragraph{GS} 中国至美国集装箱船运量连续四周环比下降，第29周TEU同比下降23\%，中国大陆降幅（-23\%）大于亚洲其他地区（-9\%），反映供应链转移趋势。洛杉矶港下周预计到港量环比增2\%，但后续一周可能转降10\%。
{\small\color{refgray}数据：第29周中国至美TEU同比-23\%\par}
{\small\color{refgray}数据：中国大陆TEU同比-23\%，亚洲其他地区-9\%\par}
{\small\color{refgray}数据：洛杉矶港下周到港量环比+2\%，后续一周同比-10\%\par}
{\small\color{refgray}边际变化：此前市场关注关税提前出货效应，GS高频数据显示该效应已消化，后续需观察7-8月补库节奏及有效税率变化对进口的影响。\par}
\paragraph{MS} AI算力需求长期供不应求，近期板块回调属技术性调整。企业token支出上限担忧缺乏数据支撑：当前月均token支出不足11美元，而AI节省劳动力成本约55美元/任务。中国大模型进展强化而非削弱算力需求，因杰文斯悖论推动总消耗增长。
{\small\color{refgray}数据：企业用户月均token支出不足11美元\par}
{\small\color{refgray}数据：AI节省劳动力成本约55美元/任务\par}
{\small\color{refgray}数据：前沿模型在超50\%人类知识任务中达到或超越人类水平\par}
{\small\color{refgray}边际变化：此前市场担忧企业限制token支出和中国模型突破削弱美国算力需求，MS指出两者均不成立，并强调AI能力非线性提升速度未被充分认知。\par}
\paragraph{MS} HDD行业处于结构性供应短缺，每年300-400EB缺口持续至2028年。2027-2028年合约价谈判已上移至25-30美元/TB，现货溢价30-40\%以上。市场对产能过剩的担忧被高估，HDD相对eSSD仍保持17-20倍成本优势。
{\small\color{refgray}数据：每年300-400EB供应缺口至2028年\par}
{\small\color{refgray}数据：合约价上移至25-30美元/TB\par}
{\small\color{refgray}数据：现货溢价30-40\%以上\par}
{\small\color{refgray}数据：HDD相对eSSD成本优势17-20倍\par}
{\small\color{refgray}边际变化：此前市场担忧HDD产能扩张和NAND价格影响，MS产业链调研显示定价谈判和客户行为均指向供应短缺持续，且HAMR产能专门用于高容量产品，不会导致总量过剩。\par}
\paragraph{NOM} 宁德时代2Q26净利润225亿元，同比+36\%，符合预期。上半年电池出货量约430-440GWh，同比+60\%，储能业务收入增速88\%领先。全年出货量预计约1TWh，2027-2028年维持19-22\%增长。毛利率环比降1.7pp至23.2\%，受原材料成本上涨影响。
{\small\color{refgray}数据：2Q26净利润225亿元，同比+36\%\par}
{\small\color{refgray}数据：上半年电池出货量约430-440GWh，同比+60\%\par}
{\small\color{refgray}数据：储能收入同比+88\%，EV电池收入同比+46\%\par}
{\small\color{refgray}数据：毛利率23.2\%，环比-1.7pp\par}
{\small\color{refgray}边际变化：此前市场关注毛利率下行压力，NOM指出出货量增长强劲且成本可控，回购计划（200-400亿元）反映公司信心，上调未来三年盈利预测。\par}
\paragraph{Bernstein} 全球储能产业正在经历一轮由技术路线、地缘政策和资本流向共同驱动的结构性调整。Bernstein在7月27日发布的全球储能周报中，梳理了北美、亚洲、欧洲及其他地区在过去一周的关键事件。这些事件单独看是公司或国家的个别动作，但放在一起，指向一个更清晰的趋势：储能产业正从单一的成本竞争，转向技术专利、供应链合规和区域化布局的多维博弈。
{\small\color{refgray}数据：报告显示，北美市场的竞争焦点正在从产能规模转向技术专利和供应链合规。LG Energy Solution向美国国际贸易委员会和联邦法院起诉中国电池制造商EVE Energy，涉及圆柱电池和无极耳电池设计等五项专利，试图阻止相关产品在美进口和销售。与此同时，韩国电池材料公司Pino正筹备高达1.6万亿韩元的融资，用于扩大磷酸铁锂正极材料产能，并降低中国股东CN...\par}
{\small\color{refgray}数据：亚洲市场的动向则更多体现在产能扩张和技术储备上。中国发布了新的可再生能源发展计划，目标到2030年将可再生能源消费量提升53\%，其中风电和光伏发电量目标为4万亿千瓦时，储能支撑的可再生能源需满足20\%的峰值电力需求。这一目标意味着储能系统在中国电力系统中的角色将从辅助转向主力。与此同时，韩国企业也在加速下一代电池技术的商业化：EcoPro BM的硫化物固态电...\par}
\paragraph{Bernstein} 一份来自Bernstein的研究报告，试图回答一个长期困扰机构读者的问题：量化模型与基本面研究，哪一种选股方法更有效？该机构的结论是，两者结合的效果显著优于单一方法。报告以2026年7月为观察节点，筛选出17只在量化排名和基本面评级上同时靠前的亚洲公司，并回溯了这一组合策略的历史表现。
{\small\color{refgray}数据：报告的核心证据来自其构建的“量化+基本面”组合的历史回测。自2013年以来，同时进入BernsteinAlpha模型最优两个五分位、且被该机构分析师给予“跑赢”评级的公司，年化超额回报达到9.4\%。如果将条件收紧至仅取最优五分位，年化超额回报进一步升至13.2\%。这一结果并非理论推演——该机构自2020年起每半年发布一次实盘组合，截至2026年6月，全部12...\par}
{\small\color{refgray}数据：KC评论：** 13.2\%的年化超额回报在量化策略回测中并不罕见，但该报告的可信之处在于，它提供了自2020年以来连续12期实盘组合的跟踪记录。读者可以关注的是，这一策略在市值加权口径下仍有两次未能跑赢基准，说明策略的稳定性并非相对，但长期胜率确实较高。\par}
\paragraph{GS} 2024年7月第三至第四周，GS分析师在亚洲（香港、新加坡）与约55名读者进行了为期八天的交流。此时正值市场对AI需求和竞争格局前景担忧加剧、市场波动性较高的阶段。但报告显示，多数读者对半导体及AI相关需求和基本面仍持正面态度，尤其对电子材料领域的兴趣保持高位。读者关注的核心问题，是相关公司能否在AI材料供需趋紧的环境下及时进行资本开支并行使定价权。讨论热度...
{\small\color{refgray}数据：读者对日东纺的悲观情绪集中体现在T-glass竞争预期上\par}
{\small\color{refgray}数据：报告指出，过去三个月左右，许多读者因日东纺主力产品T-glass供需平衡出现变化以及相对较高的估值，选择做空或低配该股。市场普遍认为，公司在低热膨胀玻璃和低介电玻璃领域面临来自台湾、日本及中国大陆新进入者的竞争，其竞争优势可能被侵蚀。然而，GS近期将日东纺评级从中性上调至观点偏积极，并认为尽管有新进入者，但竞争对手产品的实际采用进度显著延迟，公司竞争优势短期...\par}
\paragraph{GS} 沃达丰2027财年第一季度交出了一份温和超预期的成绩单。集团将全年EBITDA和自由现金流指引上调至此前区间上限，主要驱动力来自非洲业务——Vodacom在同期上调了自身指引。但欧洲有机服务收入增长仍停留在0.9\%的窄幅区间，德国和英国的增长背后都有一次性或基数因素支撑，并非结构性改善。
{\small\color{refgray}数据：这份财报的核心信号是：非洲的宏观改善和能源对冲收益帮助集团整体数字变得更好看，但欧洲的基本面增长动能并未发生变化。公司管理层对各地增长驱动因素的表述与上一季度基本一致。\par}
{\small\color{refgray}数据：Vodacom一季度有机服务收入增长12.6\%，高于市场预期的11.5\%。GS分析指出，这背后既有好于预期的宏观环境和有利的能源对冲等周期性因素，也有结构性积极因素——近期提价在整个区域落地情况良好。非洲业务的改善直接推动集团将EBITDA和自由现金流指引从此前的区间调整为区间上限，对应2027财年分别约7\%和20\%的同比增长。\par}
\subsubsection*{关键数据}
\begin{itemize}
  \item 中国广义财政赤字率从11.0\%收窄至7.0\%
  \item 土地出让收入6月同比-42\%
  \item 中国对非美发达市场出口增长约20\%
  \item 非美发达市场商品价格平均降0.7\%
  \item 1H26中国新增智能算力595 EFLOPS
  \item 2025-2028年中国数据中心CAGR 20\%
  \item HDD合约价上移至25-30美元/TB
  \item 宁德时代2Q26净利润225亿元，同比+36\%
\end{itemize}
\begin{figure}[htbp]
\centering
\includegraphics[width=0.88\linewidth]{figures/fig_001.jpg}
\caption{EXHIBIT 1 - EXHIBIT 2: On absolute basis, our last Q+F portfolio for 1H26, generated 50.4\% eq-wtd returns and 140.7\% cap-wtd returns between Jan 2026-Jun 2026}
\end{figure}
\begin{figure}[htbp]
\centering
\includegraphics[width=0.88\linewidth]{figures/fig_019.jpg}
\caption{Exhibit 1 - Exhibit 3: China's \textbackslash{}\textasciitilde{}600 EFLOPS intelligent computational power addition in 1H26 tracks below 2H25 run rate, yet we expect the pace to re-accelerate in 2H26E}
\end{figure}
\begin{figure}[htbp]
\centering
\includegraphics[width=0.88\linewidth]{figures/fig_029.jpg}
\caption{Exhibit 1 - Exhibit 2: Chinese Import Growth Has Undershot Its Pre-Pandemic Trend}
\end{figure}
\begin{figure}[htbp]
\centering
\includegraphics[width=0.88\linewidth]{figures/fig_037.jpg}
\caption{Exhibit 1 - Exhibit 1: Re-escalation of the War with Iran Has Pushed Oil and Refined Product Prices Higher Again We expect the FOMC to leave the fed funds rate unchanged at its July meeting next week. President Logan has expressed support fo}
\end{figure}
{\small\color{refgray}本节综合 12 篇研究；涉及 Bernstein、GS、MS、NOM。}
\newpage
\subsection{人民币中间价模型信号：美元走弱驱动，政策窗口临近}
\textbf{NOM模型显示人民币中间价预测值下移254个基点至6.7685，主要受美元走弱驱动，而非中国央行主动引导；逆周期因子计入后仍偏强，后续政策会议密集，中间价进入关键观察期。}
\subsubsection*{主流共识}
\begin{itemize}
  \item 美元走弱是近期人民币中间价预测下移的主要外部驱动因素
  \item 模型预测值低于当前市场收盘价，暗示中间价存在一定低估空间
  \item 逆周期因子计入后预测值仍偏强，显示政策端对汇率稳定的关注
\end{itemize}
\subsubsection*{分歧与条件差异}
\begin{itemize}
  \item 模型输出与政策意图的区分：NOM模型反映全球美元定价，不必然代表央行主动引导升值
  \item 逆周期因子调整幅度是否足以抵消市场压力，市场存在分歧
  \item 后续政策会议（如经济工作会议、货币政策报告）对中间价方向的影响方向不确定
\end{itemize}
\subsubsection*{机构观点}
\paragraph{NOM} NOM亚洲外汇策略团队更新USD/CNY中间价预测模型，不含逆周期因子的预测值为6.7685，较前次6.7939下移254个基点，较当日官方收盘价低59个基点。计入逆周期因子后预测值为6.7741，较前次中间价低198个基点。模型拆解显示，美元指数走弱是本次预测下移的核心驱动力，贡献最大负向权重，欧元、日元、英镑紧随其后。模型误差序列波动收窄，预测稳定性有所提升。报告提示后续关注7月下旬经济工作会议、8月央行二季度货币政策报告、9月外事访问及货币政策委员会会议等政策窗口。
{\small\color{refgray}数据：不含逆周期因子预测值6.7685，较前次下移254个基点\par}
{\small\color{refgray}数据：较当日官方收盘价低59个基点\par}
{\small\color{refgray}数据：计入逆周期因子后预测值6.7741，较前次中间价低198个基点\par}
{\small\color{refgray}数据：美元指数贡献最大负向权重，是预测下移核心驱动\par}
{\small\color{refgray}边际变化：相比前次预测，本次模型预测值下移幅度扩大（254 vs 前次未明确），但逆周期因子调整幅度收窄，显示模型认为中间价低估空间在缩小。\par}
\subsubsection*{关键数据}
\begin{itemize}
  \item NOM模型不含逆周期因子预测值6.7685，较前次下移254个基点
  \item 较当日官方收盘价低59个基点
  \item 计入逆周期因子后预测值6.7741，较前次中间价低198个基点
  \item 美元指数贡献最大负向权重，是预测下移核心驱动
  \item 后续关注7月下旬经济工作会议、8月央行二季度货币政策报告、9月外事访问及货币政策委员会会议
\end{itemize}
{\small\color{refgray}本节综合 1 篇研究；涉及 NOM。}
\newpage
\section{辅助信号｜战略咨询}
本板块整合波士顿咨询两份报告，聚焦两个主题：AI优先的组织设计如何从试点走向规模化价值，以及韩国股市重估背后的结构性分化与政策驱动。
\subsection{AI优先的组织设计：从试点到规模化价值}
\textbf{多数企业AI试点效果有限，根源在于组织模式仍为人类劳动力设计。真正的突破需要围绕AI智能体重构工作结构、治理与流程，而非将AI插入现有系统。}
\begin{itemize}
  \item 波士顿咨询研究显示，仅5\%的企业从AI中获得可规模化价值，核心障碍是运营模式与AI不兼容。
  \item 欧洲能源企业案例：通过重新设计端到端客户旅程，AI驱动满意度匹配人类表现，对外部服务商依赖度降低90\%以上，释放现金流用于加速转型。
  \item AI优先组织的关键转变：人类定义目标与约束，AI智能体在参数范围内自主优化路径，工作围绕动态协调而非固定角色展开。
  \item 转型节奏采用“先验证、再扩展”：优先改造高价值客户旅程，前三个月即产生运营成本节省，实现自筹资金推进。
\end{itemize}
\subsection{韩国股市重估：结构性分化与政策驱动}
\textbf{2025年KOSPI TSR约76\%，但上涨高度集中于半导体、造船、国防和核电四个行业。超过60\%上市公司仍以低于账面价值交易，下一阶段重估取决于企业能否将战略转向价值最大化。}
\begin{itemize}
  \item 2025年KOSPI总股东回报率约76\%，其中75个百分点来自资本利得，股息仅贡献1个百分点。
  \item 剔除四个关键行业后，2026年预期市净率从1.9倍降至1.0倍；半导体行业ROE预期达60\%，市值占比从29\%跃升至52\%。
  \item 政策改革（商法修订、价值提升计划、股息分离课税）与行业盈利复苏共同推动第一轮重估，但价值提升披露参与率仅约28\%。
  \item 报告指出，决定估值的是盈利质量而非规模，下一阶段重估取决于企业能否将战略重心转向企业价值最大化，并调整资本配置与股东回报。
\end{itemize}
\begin{figure}[htbp]
\centering
\includegraphics[width=0.88\linewidth]{figures/fig_075.jpg}
\caption{Exhibit 1 - KOSPI: \#1 Index Return Globally In 2025, the Korean stock market delivered an overwhelmingly strong performance relative to major global markets. The KOSPI's TSR was approximately 76\%, significantly exceeding the average of the top 10 global indices (approxima}
\end{figure}
\begin{figure}[htbp]
\centering
\includegraphics[width=0.88\linewidth]{figures/fig_078.jpg}
\caption{Exhibit 3 - \#\# 1. Market Re-rating Driven by Earnings Growth in Four Key Sectors: Semiconductors, Shipbuilding, Defense, and Nuclear Power The earnings improvement in the four key sectors drove the market's re-rating. The total KOSPI market cap expanded approximately thre}
\end{figure}
\newpage
\section{辅助信号｜智库/国际机构}
本板块聚焦多边开发银行与IMF最新研究，涵盖气候融资落地、希腊外部失衡与住房悖论、非洲财政改革及创业扶持机制，揭示结构性矛盾与政策执行层面的关键边际变化。
\subsection{气候融资进入执行阶段：ADB提前达标与项目管道成型}
\textbf{亚洲开发银行2025年气候融资承诺达135亿美元（占比51\%），提前实现2030年目标，标志多边开发银行气候融资从承诺转向可追踪资产配置。}
\begin{itemize}
  \item ADB 2025年气候融资承诺135亿美元，占年度总融资51\%，较2024年111亿美元增长显著，其中减缓88亿美元、适应45亿美元；2019-2025年累计达554亿美元。
  \item 自2015年绿色债券计划启动以来，已筹集约150亿美元，覆盖19种货币；蓝色债券自2021年纳入框架后发行约6.12亿美元，以澳元、新西兰元和瑞典克朗计价。
  \item COP30上启动的“冰川到农场”（G2F）区域项目动员35亿美元总融资，包括GCF 2.5亿美元赠款，覆盖中亚与西亚气候适应，设计三层结构：提升政府规划能力、推动可融资项目开发、试点成果导向韧性债券。
  \item 边际变化：ADB提前实现2030年气候融资占比50\%目标，项目筛选标准与资金分配机制已成熟，可供其他多边机构参考；G2F项目将气候适应转化为可操作、可融资的项目管道。
\end{itemize}
\begin{figure}[htbp]
\centering
\includegraphics[width=0.88\linewidth]{figures/fig_062.jpg}
\caption{Figure 1 - Greece's external sector has undergone significant transformation over the past decades. Pre-crisis imbalances—including large fiscal deficits, household overinvestment, and elevated labor costs—have largely been corrected. Nevertheless, the current account de}
\end{figure}
\subsection{结构性失衡与住房悖论：希腊外部竞争力与内部可负担性的双重挑战}
\textbf{IMF国别报告揭示希腊经常账户赤字重扩源于私人储蓄不足与出口附加值偏低，同时住房可负担性悖论由存量利用效率低下与收入流动性不足驱动，而非总量短缺。}
\begin{itemize}
  \item 希腊经常账户赤字居欧元区第二，净国际头寸最差；赤字主因从公共部门转向私人部门——非金融企业储蓄-投资平衡转负，家庭储蓄率低位，实际工资增长有限且生活成本挤压储蓄能力。
  \item 投入产出分析显示每1欧元出口仅产生约0.43欧元国内增加值，大量进口漏损削弱投资对国内收入循环的拉动；NGEU项目、疫后补库存及ICT/能源投资推动企业投资上升，但国内储蓄未能跟上。
  \item 住房可负担性悖论：超60\%家庭为全产权拥有者，人均住房存量居欧洲前列，但约四成家庭住房支出超可支配收入40\%；核心驱动因素为收入能力不足（资产充裕但收入流动性差），而非房贷压力。
  \item Logit模型回归显示，低收入全产权家庭住房成本负担概率显著高于高收入群体；2011-2021年家庭财富与收入联合分布变化，更多家庭持有流动性差住房资产但当期收入有限。
  \item 边际变化：报告将经常账户变动拆解为国内需求拉动、贸易条件冲击和利率传导三个渠道，为区分周期性波动与结构性变化提供可复用分析框架；住房悖论揭示资产充裕但收入不足群体的脆弱性。
\end{itemize}
\begin{figure}[htbp]
\centering
\includegraphics[width=0.88\linewidth]{figures/fig_062.jpg}
\caption{Figure 1 - Greece's external sector has undergone significant transformation over the past decades. Pre-crisis imbalances—including large fiscal deficits, household overinvestment, and elevated labor costs—have largely been corrected. Nevertheless, the current account de}
\end{figure}
\begin{figure}[htbp]
\centering
\includegraphics[width=0.88\linewidth]{figures/fig_067.jpg}
\caption{Figure 1 - Like many other European countries, Greece faces growing affordability challenges. House prices and—more recently rents—have outpaced income gains, reflecting strong and increasingly concentrated demand, including from foreigners, alongside structural supply r}
\end{figure}
\subsection{制度执行与资金配置效率：非洲财政改革与创业扶持的实证启示}
\textbf{IMF与世界银行研究分别表明，发展中国家财政收入缺口与创业扶持效果的核心瓶颈在于制度执行质量与筛选机制，而非资源总量。}
\begin{itemize}
  \item 冈比亚财产税收入仅占GDP 0.3\%，国际可比国家通常超1\%；IMF估算通过改进税制设计可额外征收约0.75\% GDP收入，缺口源于税基覆盖不全、估值体系自2005年未更新、征收手段薄弱。
  \item 冈比亚财产税征收设计缺陷：税务责任与物业所有者绑定而非物业本身，导致交易后责任追踪困难；新开发公寓楼和混合用途物业缺乏明确估值规则，大量新增房产游离于税基之外。
  \item 肯尼亚商业计划竞赛随机实验：多阶段筛选组9000美元赠款三年后使企业平均增加2.0名员工，而简化筛选组同额赠款仅增加0.6名员工（统计不显著）；筛选质量比资金规模更重要。
  \item 多阶段筛选组中36000美元赠款未带来与9000美元成比例的长期回报，表明资金规模超过阈值后边际效益递减；简化筛选虽降低行政成本，但缺乏专业评委深度评估导致资金配置效率大幅下降。
  \item 边际变化：冈比亚案例显示“把该收的收上来”依赖估值能力、登记覆盖和征收手段同步改进，而非简单加税；肯尼亚实验为高增长创业扶持提供直接证据，强调筛选机制是决定企业持续增长的核心变量。
\end{itemize}
\begin{figure}[htbp]
\centering
\includegraphics[width=0.88\linewidth]{figures/fig_072.jpg}
\caption{Figure 1 - Figure 1: Impact Evaluation Process Figure 2: Quantile Treatment Effects (3-year follow-up survey) Panel A: Number of Employees Panel B: Profits}
\end{figure}
\section{结语}
后续需验证的数据与叙事节点包括：Moderna黑色素瘤三期中期数据对肿瘤疫苗平台价值的验证；美联储7月利率路径受通胀改善与地缘政治风险的双重影响；中国数据中心全年需求20\%复合增长率的兑现情况；HDD行业合约价上移至25-30美元/TB后供应缺口能否持续至2028年；以及人民币中间价在后续经济工作会议与货币政策报告中的方向选择。
\newpage
\section{覆盖概览}
投行/券商 32 篇；战略咨询 2 篇；智库/国际机构 5 篇。\par
投行覆盖：GS 15篇 / NOM 8篇 / Bernstein 6篇 / MS 3篇\par
\section{Disclaimer}
{\scriptsize\color{refgray}
This community is a paid learning and information-sharing community created for general educational purposes only. The membership fee is charged solely for access to community discussions, educational content, organization, and operational costs. It is not an advisory fee, management fee, performance fee, brokerage fee, or compensation for personalized financial, investment, legal, tax, or professional advice.\par
I participate and share information solely as an individual and for educational discussion among community members. I am not acting as a financial adviser, investment adviser, broker, dealer, portfolio manager, legal adviser, tax adviser, accountant, fiduciary, or any other licensed professional adviser. Nothing shared in this community constitutes, and should not be interpreted as, financial advice, investment advice, legal advice, tax advice, a recommendation, solicitation, offer, endorsement, instruction, or guarantee to buy, sell, hold, short, trade, or invest in any security, cryptocurrency, fund, derivative, strategy, or other financial product.\par
All content shared in this community, including messages, discussions, links, files, charts, examples, market commentary, watchlists, AI-generated or AI-polished summaries, and third-party materials, is general, impersonal, and educational in nature. It does not take into account any individual's financial situation, investment objectives, risk tolerance, investment horizon, liquidity needs, tax status, legal restrictions, or suitability requirements. No content should be relied upon as suitable for any specific person, account, portfolio, or financial decision.\par
Information shared in this community may be incomplete, inaccurate, outdated, speculative, AI-generated, AI-edited, or based on third-party internet sources that have not been independently verified. I make no representation or warranty as to the accuracy, completeness, timeliness, reliability, or suitability of any information shared.\par
Financial markets involve substantial risk, including the possible loss of principal. Past performance, historical data, hypothetical examples, simulated results, backtests, personal opinions, or educational examples do not guarantee future results. Each member is solely responsible for conducting their own research, exercising independent judgment, and making their own decisions. Before making any financial, investment, legal, or tax decision, members should consult qualified and licensed professionals.\par
Participation in this community, payment of any membership fee, or communication with me or other members does not create any adviser-client, fiduciary, professional, contractual, agency, partnership, employment, or other advisory relationship. By joining or participating in this community, you acknowledge and agree that you are solely responsible for your own decisions, actions, transactions, profits, losses, risks, and consequences, and that I shall not be liable for any direct, indirect, incidental, consequential, financial, legal, tax, or other loss or damage arising from or related to any content, discussion, information, or material shared in this community.\par
}
\newpage
\begin{center}
{\Large 更多详情报告kcdesk.com}\par\vspace{0.8cm}
\includegraphics[width=0.58\linewidth]{\detokenize{../../prompts/zsxq_img.jpg}}
\end{center}
\end{document}
